\documentclass[11pt]{article}
\usepackage[margin=1in]{geometry}
\usepackage[T1]{fontenc}
\usepackage[utf8]{inputenc}
\usepackage{amsmath,amssymb}
\usepackage{xcolor}
\usepackage{microtype}
\usepackage[colorlinks=true,urlcolor=blue!55!black]{hyperref}

\title{The graph-topology Fredholm problem has a full K-theoretic answer}
\author{Literature-status note for arXiv:math/0401411}
\date{11 August 2026}

\begin{document}
\maketitle

\begin{center}
\fbox{\parbox{0.91\textwidth}{\textbf{Status: literature\_already\_answered.}
Let $\mathcal{CF}^{\mathrm{sa}}$ be the graph-topology space of closed,
dense-domain, self-adjoint Fredholm operators on an infinite-dimensional
separable complex Hilbert space.  Prokhorova's Theorem~C identifies this
space, by natural homotopy equivalences, with classical models for odd
$K$-theory.  Hence
\[
  \mathcal{CF}^{\mathrm{sa}}\text{ classifies }K^1,
  \qquad \pi_1(\mathcal{CF}^{\mathrm{sa}})\cong\mathbb Z.
\]
The spectral-flow map on $\pi_1$ is therefore an isomorphism, up to sign.}}
\end{center}

\section*{The original problem}

Lesch considers the set $\mathcal{CF}^{\mathrm{sa}}$ of unbounded
self-adjoint Fredholm operators with the graph metric $d_G$.  The source
paper establishes that this space is path connected and that spectral flow
gives a surjective homomorphism
\[
  \operatorname{sf}:\pi_1(\mathcal{CF}^{\mathrm{sa}},d_G)\longrightarrow
  \mathbb Z.
\]
It then asks explicitly: determine this fundamental group, and decide whether
the space is a classifying space for the $K^1$ functor
\cite[p.~3]{source}.

\section*{The homotopy-equivalence theorem}

Write ${}^{g}\!\mathcal R_F^{\mathrm{sa}}(H)$ for the space of regular
self-adjoint Fredholm operators equipped with the graph topology.  This is
the same operator space and topology as in Lesch's question.  Prokhorova
places it in a commutative diagram with the nontrivial component of bounded
self-adjoint Fredholm operators, its Riesz-topology counterpart, and the
standard Fredholm-unitary model.  Her Theorem~C states that every map in this
diagram is a homotopy equivalence and concludes that every space in the
diagram classifies $K^1$ \cite[pp.~5--6, Theorem~C]{answer}.

In particular, the theorem answers the stronger half of Lesch's problem
affirmatively:
\[
  (\mathcal{CF}^{\mathrm{sa}},d_G)\simeq U_F,
\]
where $U_F$ is a classical representing space for odd complex $K$-theory.
The paper also explicitly notes that the graph-topology space is path
connected \cite[p.~5]{answer}.

\section*{Fundamental group and spectral flow}

Since the space represents odd complex $K$-theory, Bott periodicity gives
\[
  \pi_1(\mathcal{CF}^{\mathrm{sa}},d_G)
  \cong \pi_1(U_F)
  \cong K^1(S^1)
  \cong \mathbb Z.
\]
Lesch had already proved that spectral flow is surjective onto $\mathbb Z$.
After identifying the domain with $\mathbb Z$, this surjective homomorphism is
multiplication by $+1$ or $-1$.  Thus spectral flow is an isomorphism; the
sign depends only on the normalization convention.

The result is not new to this run.  Indeed, Prokhorova's introduction records
that Joachim proved the $K^1$ analogue for regular self-adjoint operators in
2003, and presents her theorem as a transparent proof with explicit natural
homotopy equivalences \cite[pp.~2--3]{answer}.  The modern theorem directly
settles both clauses of the source problem and also makes the homotopy model
fully explicit.

\section*{Scope and verification}

The exact problem statement was checked in the official arXiv source and in
the included local compilation.  The answering statement was checked in the
official PDF of arXiv:2110.14359v3, especially Theorem~C and its introductory
setup.  The run's registry, solution, attempt, and proof-gap indexes contained
no prior resolution of this exact target.

\begin{thebibliography}{2}
\bibitem{source}
Matthias Lesch,
\emph{The uniqueness of the spectral flow on spaces of unbounded
self-adjoint Fredholm operators},
arXiv:math/0401411; in \emph{Spectral Geometry of Manifolds with Boundary and
Decomposition of Manifolds}, Contemp. Math. \textbf{366} (2005), 193--224.

\bibitem{answer}
Marina Prokhorova,
\emph{Spaces of unbounded Fredholm operators. I. Homotopy equivalences},
arXiv:2110.14359v3, 2025.
\end{thebibliography}

\end{document}
